\documentclass[12pt]{article}

\usepackage[utf8]{inputenc}
\usepackage[T2A]{fontenc}
\usepackage[serbianc]{babel}
\usepackage{indentfirst}
\usepackage{amsmath, amsthm}
\usepackage{graphicx}
\usepackage{xcolor}
\usepackage{float}
\usepackage{titling}

\title{\textbf{Кратак увод у вероватноћу}}
\author{Ђорђе Николић}
\date{\today}

\theoremstyle{definition}
\newtheorem{definition}{Дефиниција}
\begin{document}

\setlength{\droptitle}{-2.0cm}

\maketitle

\section{Увод}
Многи процеси у стварном свету су неизвесни, зато постоји потреба да се 
ове неизвесности формално опишу на математички прецизан начин.

\textbf{Случајна величина} је број који се добија као резултат неког случајног експеримента.
Да бисмо знали све о случајној величини, морамо знати које вредности она може да узме и са којом вероватноћом она узима те вредности.

Управо то представља \textbf{случајна расподела} - правило које нам говори колико је вероватно да ће исход бити овај или онај.

Случајне величине могу бити \textbf{дискретног} и \textbf{апсолутно-непрекидног} типа, али ћемо овде само навести карактеристике дискретних случајних величина.

\section{Основни појмови и дефиниције}

\begin{definition}
\textbf{Вероватноћа} догађаја \(X\), означена са \(P(X)\), је број између 0 и 1 који квантитативно изражава колико је вероватно да ће се тај догађај десити. 
\begin{itemize}
    \item \(P(X) = 0\) означава немогућ догађај.
    \item \(P(X) = 1\) означава сигуран догађај.
\end{itemize}

\end{definition}
\begin{definition}
Случајна величина \(X\) је \textbf{дискретног типа} ако узима коначан или пребројив скуп вредности и важи:
\[
P(X = x_i) = p_i, \quad i = 1, 2, \dots
\]
где је \(p_i \ge 0\) и
\[
\sum_{i} p_i = 1.
\]

\end{definition}

\section{Неке познате случајне расподеле}
\begin{itemize}
\item Бернулијева расподела
\item Биномна расподела
\item Геометријска расподела
\item Униформна расподела

\end{itemize}

\begin{table}[H]
\centering
\renewcommand{\arraystretch}{1.5}
\begin{tabular}{|c|c|}
\hline
\textbf{Тип расподеле} & \textbf{Закон расподеле} \\ \hline
Бернулијева $X \sim \text{Bern}(p)$ & $P(X=1)=p, \quad P(X=0)=1-p$   \\ \hline
Биномна $X \sim \text{Bin}(n,p)$ &  $P(X = k) = \binom{n}{k} p^k (1-p)^{\,n-k}, \quad k= 0,1,\dots,n$ \\ \hline
Геометријска $X \sim \text{Geom}(p)$ & $P(X=k) = (1-p)^{k-1} p, \quad k=1,2,3,\dots$  \\ \hline
Униформна $X \sim \text{Uniform}\{1,2,\dots,n\}$ & $P(X = k) = \frac{1}{n} , \quad k =1,\dots , n$  \\ \hline
\end{tabular}
\caption{Примери неких расподела}
\end{table}

\subsection{Бернулијева расподела}

Бернулијева расподела је дискретна расподела са два исхода:
\emph{успех} и \emph{неуспех}.  
Вероватноћа успеха износи $p$, док је вероватноћа неуспеха $1-p$.

\textbf{Пример.} При бацању новчића:
глава представља успех, а писмо неуспех.

\subsection{Биномна расподела}

Биномна расподела је дискретна расподела која описује број успеха у $n$ независних понављања Бернулијевог огледа.
Вероватноћа успеха у сваком понављању је $p$, док је вероватноћа неуспеха $1-p$.


\textbf{Пример.} Бацање новчића $n$ пута и бројање колико се пута појавила глава.


\begin{figure}[H]
    \centering
    \includegraphics[width=0.7\textwidth]{Binom.png}
    \caption{Граф $\text{Bin}(20,0.5)$ расподеле}
    \label{fig:биномна}
\end{figure}

\subsection{Геометријска расподела}

Геометријска расподела је дискретна расподела која описује број понављања Бернулијевог огледа потребних да би се први пут појавио успех.
Вероватноћа успеха у сваком понављању је $p$.

\textbf{Пример.} Бацање новчића све док се први пут не појави глава.
\subsection{Униформна расподела}

Униформна расподела је расподела код које су сви могући исходи подједнако вероватни.
Може бити дискретна или апсолутно-непрекидна, у зависности од скупа могућих исхода.

\textbf{Пример.} Насумичан избор једног броја из скупа $\{1,2,3,4,5,6\}$ при бацању правилне коцке.


\section{Закључак}
Вероватноћа представља важну област математике која омогућава описивање и анализу случајних појава. Дискретне расподеле, као што су Бернулијева, Биномна, Геометријска и Униформна, имају широку примену у статистици, информатици и многим другим наукама. Разумевање ових појмова представља добру основу за даље проучавање статистике и случајних процеса.

\end{document}
